Electromagnetic Dissociation (EMD) is a process in which the nucleus or nucleon is excited by a photon, and emits nucleons as it comes down to the ground state. This effect relies on high-energy photon impinging upon the target nucleus, and therefore its' cross section increases with $\sqrt{s_{NN}}$. The exact process that takes place depends on the incoming photon energy, as can be seen in Fig. \ref{fig:emd_processes_photon_energy}. For heavy nuclei such as Au and Pb, up to ($\sim$10 GeV) the dominant contribution is known to be from the Giant Dipole Resonance (GDR), which is a process in which the excitation of the nucleus causes the protons to move in unison with respect to the neutrons, generating a dipole within the nucleus itself. De-excitation of GDR is usually accompanied by neutron emission, and it is one of the mechanisms that generate spectator neutrons which are read by the Zero Degree Calorimeter (see section \ref{sec:atlas}). 

\begin{figure}[h!] 
	\centering % Centers the image and caption
	\includegraphics[width=0.7\textwidth]{\main/illustrations/pshenichnov_EMD_regions.png}
	\caption{Product of the spectrum of equivalent photons and total photonuclear cross section for 10.6 A GeV gold nuclei (dotted curve) and 158 A GeV lead nuclei (solid curve) in the interactions with gold and lead targets, respectively. The regions of GR excitation, quasi-deuteron absorption,and $\Delta$-isobar excitation are shown. For comparison, a product of the photon spectrum at an energy of 158 A GeV and the partial cross sections of photonuclear reactions with \ce{^{83}_{}Bi} (dashed curve) is shown \cite{pshenichnov_electromagnetic_2011}.}
	\label{fig:emd_processes_photon_energy}
\end{figure}

As the photon energy increases, other channels become available due to the higher resolution provided by a shorter light pulse, and the process of interaction with quasi-deuteron (np pair) is unlocked. Eventually, with photon energy of $\gtrsim 100$ MeV the channel for interaction with individual nucleons is possible. Higher probe energies will allow the photon to interact with the nucleon constituents, the quarks and gluons, and will be the topic of the next sections.
